\documentclass[11pt,a4paper]{article}

% ---------- Encodage et langue ----------
\usepackage[utf8]{inputenc}
\usepackage[T1]{fontenc}
% \usepackage{lmodern}  % décommenter si disponible (police vectorielle)
\usepackage[french]{babel}
\usepackage[a4paper,margin=2.3cm]{geometry}

% ---------- Mathématiques ----------
\usepackage{amsmath,amssymb,amsthm}
\usepackage{mathtools}
\usepackage{enumitem}
\usepackage[most]{tcolorbox}
\usepackage[colorlinks=true,linkcolor=blue!60!black,urlcolor=blue!60!black]{hyperref}

% ---------- Macros ----------
\newcommand{\R}{\mathbb{R}}
\newcommand{\N}{\mathbb{N}}
\newcommand{\Z}{\mathbb{Z}}
\newcommand{\Q}{\mathbb{Q}}
\newcommand{\C}{\mathbb{C}}
\newcommand{\K}{\mathbb{K}}
\newcommand{\Sun}{\mathbb{S}^1}
\newcommand{\interieur}[1]{\mathring{#1}}
\newcommand{\adh}[1]{\overline{#1}}
\newcommand{\Int}{\operatorname{int}}
\newcommand{\diam}{\operatorname{diam}}
\newcommand{\dist}{\operatorname{dist}}
\newcommand{\Vect}{\operatorname{Vect}}
\newcommand{\Bf}{B_f}
\newcommand{\norm}[1]{\left\|#1\right\|}
\newcommand{\abs}[1]{\left|#1\right|}
\newcommand{\opnorm}[1]{|\!|\!|#1|\!|\!|}
\newcommand{\ps}[2]{\left\langle #1,#2\right\rangle}
\newcommand{\Lc}{\mathcal{L}}
\newcommand{\Cz}{\mathcal{C}^0}
\newcommand{\Cinf}{\mathcal{C}^\infty}
\newcommand{\Cc}{\mathcal{C}^0_c}
\newcommand{\ie}{c'est-à-dire}
\newcommand{\wk}{\rightharpoonup}
\newcommand{\wks}{\rightharpoonup^{\!\star}}
\DeclareMathOperator{\Ker}{Ker}
\DeclareMathOperator{\Ima}{Im}
\DeclareMathOperator{\supp}{supp}
\DeclareMathOperator{\sgn}{sgn}
\renewcommand{\Re}{\operatorname{Re}}
\renewcommand{\Im}{\operatorname{Im}}
\newcommand{\D}{\mathcal{D}}
\newcommand{\Sha}{\mathrm{III}}

% ---------- Environnements ----------
\newtcolorbox{exo}[1]{
  enhanced, breakable,
  colback=black!3, colframe=black!70,
  fonttitle=\bfseries, title={#1},
  boxrule=0.6pt, arc=1pt, left=6pt, right=6pt
}
\newtcolorbox{pointcle}{
  enhanced, breakable,
  colback=blue!4, colframe=blue!55!black,
  fonttitle=\bfseries, title={Point clé},
  boxrule=0.8pt, arc=2pt, left=6pt, right=6pt
}
\newtcolorbox{piege}{
  enhanced, breakable,
  colback=red!4, colframe=red!60!black,
  fonttitle=\bfseries, title={Piège / difficulté},
  boxrule=0.8pt, arc=2pt, left=6pt, right=6pt
}
\newtcolorbox{commentaire}{
  enhanced, breakable,
  colback=gray!8, colframe=gray!55!black,
  fonttitle=\bfseries, title={Commentaire},
  boxrule=0.6pt, arc=2pt, left=6pt, right=6pt
}
\newtcolorbox{methode}{
  enhanced, breakable,
  colback=green!4, colframe=green!40!black,
  fonttitle=\bfseries, title={Méthode},
  boxrule=0.8pt, arc=2pt, left=6pt, right=6pt
}

\theoremstyle{definition}
\newtheorem*{solution}{Solution}

\title{\textbf{4MA305 -- Bases d'analyse fonctionnelle}\\[2mm]
\Large TD n\textsuperscript{o}6 -- Distributions (énoncés)}
\author{}
\date{Année universitaire 2026--2027}

\begin{document}
\maketitle

\begin{commentaire}
Feuille d'exercices sur le chapitre~6 du polycopié, complémentaire des exercices du poly et des développements du TD0 (chapitre~6). Notations : $\D(\Omega)$, $\D'(\Omega)$, $\langle T,\varphi\rangle$, $h=\mathbf 1_{\R_+}$, $\delta_a$, $\mathrm{vp}(1/x)$, $\check\psi(x)=\psi(-x)$, $\tau_af=f(\cdot-a)$ ; on identifie $f\in L^1_{\mathrm{loc}}$ et $T_f$. Le corrigé détaillé est dans \texttt{corrige\_TD6.pdf} ; l'idée est de chercher les exercices en conditions d'examen avant de le lire. Difficulté indicative : 1, 3, 4 (fondamentaux) ; 2, 6, 7 (classiques d'examen) ; 5, 8, 9, 10 (plus longs, à travailler après les autres).
\end{commentaire}

\section*{Fonctions tests, distributions, convergence}

\begin{exo}{Exercice 1 -- Fabriquer des fonctions tests}
\begin{enumerate}[label=(\alph*)]
  \item Soit $\gamma(t):=e^{-1/t}$ pour $t>0$ et $\gamma(t):=0$ pour $t\le0$. Montrer que $\gamma\in\Cinf(\R)$ et que $\gamma^{(k)}(0)=0$ pour tout $k$. \emph{Indication : montrer par récurrence que $\gamma^{(k)}(t)=P_k(1/t)e^{-1/t}$ pour $t>0$, avec $P_k$ un polynôme.}
  \item En déduire une fonction $\rho\in\D(\R^d)$, positive, paire, d'intégrale $1$, à support dans $\Bf(0,1)$ et strictement positive sur $B(0,1)$.
  \item (Fonction plateau en dimension 1.) Soient $a<b$ et $\varepsilon>0$. Construire explicitement $\theta\in\D(\R)$ à valeurs dans $[0,1]$, valant $1$ sur $[a,b]$ et $0$ hors de $]a-\varepsilon,b+\varepsilon[$. \emph{Indication : primitive d'une différence de deux bosses.}
  \item Montrer qu'il n'existe pas de fonction test non nulle qui soit développable en série entière au voisinage de chaque point de $\R$ (analytique). Pourquoi est-ce important pour la théorie ?
  \item Soit $\varphi\in\D(\R)$ et $\varphi_\lambda(x):=\varphi(x/\lambda)$ pour $\lambda>0$. Calculer $\norm{\varphi_\lambda}_{m,\R}$ en fonction de $\lambda$ et des normes de $\varphi$. La famille $(\varphi_\lambda)_{\lambda\to0}$ converge-t-elle dans $\D(\R)$ ? Et $(\varphi_\lambda)_{\lambda\to\infty}$ ?
\end{enumerate}
\end{exo}

\begin{exo}{Exercice 2 -- Distributions ou pas ? Ordre fini ou infini ?}
\begin{enumerate}[label=(\alph*)]
  \item Montrer que $T_1:=\sum_{k\in\Z}\delta_k$ est une distribution d'ordre $0$ sur $\R$.
  \item Montrer que $T_2:=\sum_{k\ge1}\delta_k^{(k)}$ (\ie\ $\langle T_2,\varphi\rangle=\sum_{k\ge1}(-1)^k\varphi^{(k)}(k)$) est une distribution sur $\R$, mais qu'elle n'est pas d'ordre fini.
  \item Soit $T_3(\varphi):=\sum_{k\ge1}\varphi(1/k)$. Montrer que $T_3$ définit une distribution d'ordre $0$ sur $\Omega=]0,+\infty[$, mais que la même formule ne définit pas une distribution sur $\R$ (ni même une forme linéaire sur $\D(\R)$).
  \item Montrer que $T_4(\varphi):=\sum_{k\ge1}\big(\varphi(1/k)-\varphi(0)-\tfrac1k\varphi'(0)\big)$ définit une distribution sur $\R$, d'ordre $\le2$.
\end{enumerate}
\end{exo}

\begin{exo}{Exercice 3 -- Limites au sens des distributions}
Déterminer les limites dans $\D'(\R)$, quand $n\to\infty$, des suites suivantes.
\begin{enumerate}[label=(\alph*)]
  \item $f_n:=n\mathbf 1_{[0,1/n]}$.
  \item $g_n:=n^2\big(\mathbf 1_{[0,1/n]}-\mathbf 1_{[-1/n,0]}\big)$.
  \item $S_n:=n(\delta_{1/n}-\delta_0)$.
  \item $u_n(x):=\cos^2(nx)$.
  \item $v_n(x):=\dfrac{\sin(nx)}{\pi x}$. \emph{Indication : on admettra $\int_{-\infty}^{+\infty}\frac{\sin u}u\,du=\pi$ ; séparer $\varphi(x)=\varphi(0)\chi(x)+(\varphi(x)-\varphi(0)\chi(x))$ avec $\chi$ une fonction plateau valant $1$ près de $0$, et utiliser le lemme de Riemann--Lebesgue.}
  \item Existe-t-il une suite $(\varphi_n)$ de fonctions tests convergeant vers $\delta_0$ dans $\D'(\R)$ ? convergeant vers $\delta_0$ dans $\D(\R)$ ?
\end{enumerate}
\end{exo}

\section*{Dérivation}

\begin{exo}{Exercice 4 -- Calculs de dérivées et équations différentielles}
\begin{enumerate}[label=(\alph*)]
  \item Calculer dans $\D'(\R)$ les dérivées première et seconde de $f(x)=|\sin x|$ et de $g(x)=e^{-|x|}$.
  \item En déduire que $E:=-\frac12e^{-|x|}$ vérifie $E''-E=\delta_0$ (on dit que $E$ est une \emph{solution fondamentale} de l'opérateur $\frac{d^2}{dx^2}-1$).
  \item Déterminer toutes les distributions $T\in\D'(\R)$ telles que $T''-T=\delta_0$. \emph{Indication : pour $v\in\D'(\R)$ avec $v''=v$, considérer $(e^{-x}(v'+v))'$.}
  \item Déterminer toutes les $T\in\D'(\R)$ telles que $T''+T=\delta_0$ ; vérifier que $E_1:=h(x)\sin x$ en est une.
  \item Soit $f$ continue sur $\R$, de classe $\mathcal C^1$ sur $\R\setminus\{0\}$, telle que $f'$ (définie hors de $0$) soit localement intégrable. Montrer que $T_f'=T_{f'}$. En quoi est-ce plus général que la formule des sauts ? Donner un exemple où $f'$ n'est pas bornée près de $0$.
\end{enumerate}
\end{exo}

\begin{exo}{Exercice 5 -- Primitives d'une distribution}
Soit $I\subset\R$ un intervalle ouvert et $T\in\D'(I)$.
\begin{enumerate}[label=(\alph*)]
  \item Soit $\theta\in\D(I)$ d'intégrale $1$. Montrer que pour toute $\varphi\in\D(I)$, la fonction $\psi:=\varphi-\big(\int_I\varphi\big)\theta$ admet une unique primitive $\Psi\in\D(I)$, et que $\varphi\mapsto\Psi$ est linéaire et vérifie $\norm\Psi_{m+1,K'}\le C_K\norm\varphi_{m,K'}$ pour $\varphi\in\D_K(I)$, avec $K':=\mathrm{conv}(K\cup\supp\theta)$ (enveloppe convexe : un segment de $I$) et des constantes convenables (préciser).
  \item On pose $\langle S,\varphi\rangle:=-\langle T,\Psi\rangle$. Montrer que $S\in\D'(I)$ et que $S'=T$.
  \item Montrer que les primitives de $T$ (les $U\in\D'(I)$ avec $U'=T$) sont exactement les $S+c$, $c\in\K$.
  \item Application : soit $g\in L^1_{\mathrm{loc}}(I)$ et $U\in\D'(I)$ avec $U'=T_g$. Montrer que $U$ est une fonction continue, de la forme $U=T_F$ avec $F(x)=c+\int_{x_0}^xg$.
  \item Montrer qu'une primitive d'une distribution d'ordre $m\ge1$ est d'ordre $\le m-1$, et qu'une primitive d'une distribution d'ordre $0$ est d'ordre $0$. Quel est l'ordre des primitives de $\delta_0$ ? de $\delta_0'$ ?
\end{enumerate}
\end{exo}

\section*{Multiplication, support, valeur principale}

\begin{exo}{Exercice 6 -- Équations $xT=\dots$ et distributions à support ponctuel}
On rappelle (TD0 ch.~6, lemme d'Hadamard) que si $\varphi\in\D(\R)$ vérifie $\varphi(0)=0$, alors $\varphi=x\psi$ avec $\psi\in\D(\R)$, et que les solutions de $xT=0$ sont les $c\delta_0$.
\begin{enumerate}[label=(\alph*)]
  \item Résoudre dans $\D'(\R)$ : $xT=\delta_0$ ; puis $x^2T=0$ ; puis $x^kT=0$ ($k\ge1$).
  \item Soient $a\ne b$. Résoudre $(x-a)(x-b)T=0$. \emph{Indication : écrire $T=\theta T+(1-\theta)T$ avec $\theta\in\D(\R)$ valant $1$ près de $a$ et $0$ près de $b$, et observer que $x-b$ est « inversible » sur le support de $\theta$.}
  \item On dit que $T\in\D'(\Omega)$ est \emph{nulle sur l'ouvert} $U\subset\Omega$ si $\langle T,\varphi\rangle=0$ pour toute $\varphi\in\D(U)$. Montrer, à l'aide du Corollaire~6.1, que si $T$ est nulle sur chaque ouvert d'une famille $(U_i)_{i\in I}$, elle est nulle sur $\bigcup_iU_i$. On définit le \emph{support} de $T$ comme le complémentaire du plus grand ouvert sur lequel $T$ est nulle. Déterminer les supports de $T_f$ ($f$ continue), de $\delta_a^{(k)}$, de $\mathrm{vp}(1/x)$.
  \item Soit $T\in\D'(\R)$ de support $\{0\}$ et d'ordre $\le m$ (une distribution à support compact est d'ordre fini, TD0 ch.~6). Montrer que $x^{m+1}T=0$, puis que $T=\sum_{j=0}^mc_j\delta_0^{(j)}$. \emph{Indication : pour $\varphi\in\D(\R)$, écrire $\langle T,x^{m+1}\varphi\rangle=\langle T,\chi_\varepsilon x^{m+1}\varphi\rangle$ avec $\chi_\varepsilon(x)=\chi(x/\varepsilon)$, et faire $\varepsilon\to0$.}
\end{enumerate}
\end{exo}

\begin{exo}{Exercice 7 -- Valeur principale, partie finie, formule de Sokhotski--Plemelj}
\begin{enumerate}[label=(\alph*)]
  \item On pose $\mathrm{Pf}(1/x^2):=-\big(\mathrm{vp}(1/x)\big)'$. Montrer que pour $\varphi\in\D(\R)$,
  \[
  \big\langle\mathrm{Pf}(1/x^2),\varphi\big\rangle=\lim_{\varepsilon\to0}\Big(\int_{|x|>\varepsilon}\frac{\varphi(x)}{x^2}\,dx-\frac{2\varphi(0)}\varepsilon\Big)=\int_\R\frac{\varphi(x)-\varphi(0)-x\varphi'(0)\mathbf 1_{|x|<1}}{x^2}\,dx .
  \]
  \item Montrer que $x\,\mathrm{Pf}(1/x^2)=\mathrm{vp}(1/x)$ et $x^2\,\mathrm{Pf}(1/x^2)=1$. Quelles sont toutes les solutions de $x^2T=1$ ?
  \item Résoudre dans $\D'(\R)$ l'équation $xT'+T=0$.
  \item (Sokhotski--Plemelj.) Pour $\varepsilon>0$, soit $f_\varepsilon(x):=\dfrac1{x-i\varepsilon}$. Montrer que $f_\varepsilon\in L^1_{\mathrm{loc}}(\R)$ et que
  \[
  f_\varepsilon\xrightarrow[\varepsilon\to0^+]{}\mathrm{vp}(1/x)+i\pi\delta_0\quad\text{dans }\D'(\R).
  \]
  En déduire la limite de $\frac1{x+i\varepsilon}$, et retrouver $x\,\mathrm{vp}(1/x)=1$.
  \item Montrer que $\mathrm{vp}(1/x)$ est la limite dans $\D'(\R)$ de $\dfrac{x}{x^2+\varepsilon^2}$, et que $\delta_0$ est la limite de $\dfrac1\pi\dfrac{\varepsilon}{x^2+\varepsilon^2}$ (noyau de Poisson).
\end{enumerate}
\end{exo}

\section*{Convolution, solutions fondamentales, dimension supérieure}

\begin{exo}{Exercice 8 -- Convolution et solutions fondamentales}
Soit $\psi\in\D(\R)$.
\begin{enumerate}[label=(\alph*)]
  \item Calculer $\psi\star\delta_a$, $\psi\star\delta_0^{(k)}$, $\psi\star h$ et $\psi\star\mathbf 1_{[a,b]}$ (comme fonctions lisses, Proposition~6.10).
  \item Soit $E\in\D'(\R)$ une solution fondamentale de $P(\frac d{dx})=\sum_{j\le n}a_j\frac{d^j}{dx^j}$ ($a_j\in\C$), \ie\ $P(\frac d{dx})E=\delta_0$. Montrer que $u:=\psi\star E$ est une fonction lisse vérifiant $P(\frac d{dx})u=\psi$.
  \item Avec $E=he^{-x}$ (TD0 ch.~6) puis $E=h\sin x$ (Exercice~4), expliciter la solution $u$ de $u'+u=\psi$ puis de $u''+u=\psi$ ainsi obtenue. Quelle propriété supplémentaire a cette solution parmi toutes les solutions ?
  \item Montrer que si $T\in\D'(\R)$ vérifie $T'+T=0$ alors $T=ce^{-x}$, et en déduire que $u=\psi\star(he^{-x})$ est l'unique solution de $u'+u=\psi$ nulle au voisinage de $-\infty$.
  \item Soit $T\in\D'(\R)$ et $\psi\in\D(\R)$. Montrer que $(\psi\star T)'=\psi'\star T=\psi\star T'$, puis que $\supp(\psi\star T)\subset\supp\psi+\supp T$ (somme de Minkowski). Que peut-on dire de $\psi\star T$ lorsque $\supp T$ est compact ?
\end{enumerate}
\end{exo}

\begin{exo}{Exercice 9 -- Quelques calculs en dimension $d\ge2$}
\begin{enumerate}[label=(\alph*)]
  \item Dans $\R^2$, soit $Q:=\{x_1>0,\ x_2>0\}$ et $T:=\mathbf 1_Q$. Calculer $\partial_1T$ et $\partial_1\partial_2T$.
  \item Dans $\R^2$, soit $H:=\mathbf 1_{\{x_1>0\}}$. Calculer $\partial_1H$ et $\partial_2H$, et montrer que $\partial_1H$ n'est pas une fonction.
  \item Soit $E(x):=-\dfrac1{4\pi\norm x}$ sur $\R^3$. On sait que $\Delta E=\delta_0$ (Proposition~6.5). Pour $\psi\in\D(\R^3)$ on pose $u:=\psi\star E$, \ie\ $u(x)=-\dfrac1{4\pi}\displaystyle\int_{\R^3}\frac{\psi(y)}{\norm{x-y}}\,dy$ (potentiel newtonien). Montrer que $u\in\Cinf(\R^3)$, que $\Delta u=\psi$, et que $u(x)=O(1/\norm x)$ à l'infini.
  \item En déduire que si $v\in\D'(\R^3)$ vérifie $\Delta v=\psi$, alors $v-u$ est une fonction lisse harmonique ; et que $u$ est l'unique solution de $\Delta v=\psi$ qui soit une fonction lisse tendant vers $0$ à l'infini et dont les dérivées premières sont dans $L^2(\R^3)$. \emph{On utilisera le Corollaire~6.3.}
\end{enumerate}
\end{exo}

\section*{Distributions périodiques et séries de Fourier}

\begin{exo}{Exercice 10 -- Distributions périodiques}
On dit que $T\in\D'(\R)$ est \emph{$2\pi$-périodique} si $\tau_{2\pi}T=T$. On fixe $\theta\in\D(\R)$ telle que $\sum_{n\in\Z}\theta(x+2\pi n)=1$ pour tout $x\in\R$ (une telle fonction existe : prendre $\theta:=\theta_0/\Theta_0$ avec $\theta_0$ une bosse positive valant $1$ sur $[0,2\pi]$ et $\Theta_0:=\sum_n\theta_0(\cdot+2\pi n)$, lisse, $2\pi$-périodique et $\ge1$).
\begin{enumerate}[label=(\alph*)]
  \item Montrer que $T_f$ est $2\pi$-périodique si $f\in L^1_{\mathrm{loc}}$ est $2\pi$-périodique, et que le peigne $\Sha:=\sum_{n\in\Z}\delta_{2\pi n}$ est $2\pi$-périodique. Montrer que les dérivées d'une distribution périodique sont périodiques.
  \item Soit $\psi\in\D(\R)$ avec $\sum_n\psi(x+2\pi n)=0$ pour tout $x$. Montrer que $\psi=\tau_{2\pi}\chi-\chi$ pour une $\chi\in\D(\R)$. \emph{Indication : $\chi:=\sum_{n\ge1}\psi(\cdot+2\pi n)$.} En déduire que si $T$ est $2\pi$-périodique, $\langle T,\psi g\rangle=0$ pour toute fonction $g\in\Cinf(\R)$ $2\pi$-périodique.
  \item Pour $T$ $2\pi$-périodique et $k\in\Z$, on pose $c_k(T):=\frac1{2\pi}\langle T,\theta\,e^{-ikx}\rangle$. Montrer que $c_k(T)$ ne dépend pas du choix de $\theta$, et que pour $T=T_f$ on retrouve $c_k(f)=\frac1{2\pi}\int_0^{2\pi}f(x)e^{-ikx}\,dx$. Calculer $c_k(\Sha)$.
  \item Montrer que $c_k(T')=ik\,c_k(T)$.
  \item Montrer qu'il existe $m\in\N$ et $C>0$ tels que $|c_k(T)|\le C(1+|k|)^m$ pour tout $k\in\Z$ (croissance polynomiale).
  \item En déduire que la série $\sum_{k\in\Z}c_k(T)e^{ikx}$ converge dans $\D'(\R)$ (TD0 ch.~6, séries trigonométriques). On admet que sa somme est $T$ (théorème de Fourier pour les distributions périodiques). Vérifier ce résultat sur $\Sha$ à l'aide du TD0 ch.~6 (peigne de Dirac), et sur $T=\Sha'$.
\end{enumerate}
\end{exo}

\section*{Vers l'examen : questions de cours, problème transversal, compléments}

\begin{commentaire}
Ces quatre exercices complètent la feuille dans l'optique de l'examen. L'Exercice~11 fait restituer les preuves du polycopié ; l'Exercice~12 est un problème long qui relie les chapitres 4, 5 et 6 autour de la Proposition~6.6 ; l'Exercice~13 traite deux points du polycopié absents des exercices précédents (la convexité de l'Exemple~6.11 et un opérateur non hypoelliptique, à comparer à la Remarque~6.9) ; l'Exercice~14 est un vrai ou faux sur les pièges du chapitre.
\end{commentaire}

% =====================================================================
\begin{exo}{Exercice 11 -- Questions de cours}
\emph{À traiter sans document, en rédigeant comme le jour de l'examen : énoncé complet, puis démonstration.}
\begin{enumerate}[label=(\alph*)]
  \item Énoncer et démontrer le Théorème~6.2 (une forme linéaire positive sur $\D(\Omega)$ est une distribution d'ordre $0$).
  \item Énoncer et démontrer le lemme de du Bois-Reymond (Lemme~6.2).
  \item Énoncer la formule des sauts (Théorème~6.3) et la démontrer lorsque $A=\{a\}$ est réduit à un point.
  \item Soit $I\subset\R$ un intervalle ouvert et $T\in\D'(I)$ telle que $T'=0$. Démontrer que $T$ est constante (Théorème~6.4 en dimension~$1$ ; on redémontrera le cas utile du Lemme~6.3).
  \item Énoncer et démontrer la Proposition~6.2 ($|T(\varphi)|\lesssim\norm\varphi_p$ entraîne $T\in L^{p'}(\Omega)$, pour $p<\infty$). Que se passe-t-il pour $p=\infty$ ?
\end{enumerate}
\end{exo}

% =====================================================================
\begin{exo}{Exercice 12 -- Problème : séries de Fourier, dérivée dans $L^2$ et résolution spectrale (chapitres 4, 5 et 6)}
On note $J:=\,]0,2\pi[$ et $H:=L^2(J)$ muni du produit scalaire $\langle f,g\rangle:=\frac1{2\pi}\int_Jf\bar g$, pour lequel $e_k:x\mapsto e^{ikx}$ ($k\in\Z$) est une base hilbertienne (Exemple~4.7) ; $c_k(f):=\langle f,e_k\rangle$ et $S_Nf:=\sum_{|k|\le N}c_k(f)e_k$. Pour $f\in H$ on note $f^\sharp$ son prolongement $2\pi$-périodique à $\R$ (défini presque partout), qui est dans $L^2_{\mathrm{loc}}(\R)$, et l'on considère $S_Nf$ comme une fonction sur $\R$. On pose enfin
\[
H^1_\sharp:=\Big\{u\in H:\ \sum_{k\in\Z}k^2|c_k(u)|^2<\infty\Big\},\qquad L^2_0:=\{f\in H:\ c_0(f)=0\}.
\]
\begin{enumerate}[label=(\alph*)]
  \item Soit $f\in H$. Montrer que $S_Nf\to f^\sharp$ dans $L^2([a,b])$ pour tout segment $[a,b]$, puis dans $\D'(\R)$. En déduire que $\sum_{|k|\le N}ik\,c_k(f)e_k\wk(f^\sharp)'$ dans $\D'(\R)$.
  \item Soit $u\in H^1_\sharp$. Montrer qu'il existe $v\in H$ tel que $(u^\sharp)'=v^\sharp$ dans $\D'(\R)$ ; on note $u':=v$. Montrer que la série de Fourier de $u$ converge normalement sur $\R$, que $u$ admet un représentant continu $\tilde u$ sur $[0,2\pi]$ avec $\tilde u(0)=\tilde u(2\pi)$, et que
  \[
  \norm{\tilde u-c_0(u)}_\infty\le\frac\pi{\sqrt3}\,\norm{u'}_H .
  \]
  Réciproquement, montrer à l'aide de l'Exercice~10 que si $u,w\in H$ vérifient $(u^\sharp)'=w^\sharp$ dans $\D'(\R)$, alors $u\in H^1_\sharp$.
  \item Soit $f(x):=x$ sur $J$. Calculer $c_k(f)$. Vérifier que $T_f'=T_1$ dans $\D'(J)$, mais que $f\notin H^1_\sharp$. Calculer $(f^\sharp)'$ dans $\D'(\R)$ de deux façons et commenter.
  \item (Proposition~6.6.) Soit $f\in L^2_0$. Montrer qu'il existe un unique $\varphi_f\in L^2_0$ tel que $(\varphi_f^\sharp)''=f^\sharp$ dans $\D'(\R)$, et l'expliciter. Montrer que $\varphi_f$ admet un représentant de classe $\mathcal C^1$. Montrer enfin que l'unicité est \emph{fausse} si l'on demande seulement $\varphi''=f$ dans $\D'(J)$.
  \item Montrer que $K:f\mapsto\varphi_f$ est un opérateur linéaire continu, auto-adjoint et compact de l'espace de Hilbert $L^2_0$ ; calculer $\opnorm K$, ses valeurs propres et ses espaces propres, et comparer avec la Proposition~4.7 et le Théorème~5.5.
  \item On munit $H^1_\sharp$ de $\langle u,v\rangle_1:=\langle u,v\rangle+\langle u',v'\rangle$. Montrer que c'est un espace de Hilbert séparable. Soit $f\in H$ ; à l'aide du Théorème~4.6 de Riesz, montrer qu'il existe un unique $u\in H^1_\sharp$ tel que $\langle v,u\rangle_1=\langle v,f\rangle$ pour tout $v\in H^1_\sharp$, calculer ses coefficients de Fourier et vérifier que $-(u^\sharp)''+u^\sharp=f^\sharp$ dans $\D'(\R)$.
\end{enumerate}
\end{exo}

% =====================================================================
\begin{exo}{Exercice 13 -- Différences secondes, convexité, équation des ondes}
\textbf{A. Convexité (Exemple~6.11).} Pour $T\in\D'(\R)$ et $h\ne0$ on pose $D_hT:=\dfrac{\tau_hT+\tau_{-h}T-2T}{h^2}$.
\begin{enumerate}[label=(\alph*)]
  \item Montrer que $D_hT\wk T''$ dans $\D'(\R)$ quand $h\to0$. En déduire la limite de $n^2\big(\delta_{1/n}+\delta_{-1/n}-2\delta_0\big)$.
  \item Soit $f:\R\to\R$ convexe (donc continue). Montrer que $T_f''$ est une distribution positive, donc une mesure de Radon positive.
  \item Réciproquement, soit $f:\R\to\R$ continue telle que $T_f''$ soit positive. Montrer que $f$ est convexe. \emph{Indication : régulariser par $f\star\rho_n$.}
  \item Les fonctions $|x|$, $e^{-|x|}$ et $|x^2-1|$ sont-elles convexes ? Répondre en calculant leur dérivée seconde dans $\D'(\R)$.
\end{enumerate}
\textbf{B. Équation des ondes.} On fixe $c>0$ et l'on note $\square:=\partial_{tt}-c^2\partial_{xx}$ sur $\R^2$, de variables $(x,t)$.
\begin{enumerate}[label=(\alph*),resume]
  \item Soit $f\in L^1_{\mathrm{loc}}(\R)$ et $u(x,t):=f(x-ct)$. Montrer que $u\in L^1_{\mathrm{loc}}(\R^2)$, que $\partial_tu+c\,\partial_xu=0$ dans $\D'(\R^2)$, puis que $\square u=0$. L'opérateur $\square$ vérifie-t-il un analogue du lemme de Weyl ?
  \item Soit $E:=\frac1{2c}\mathbf 1_{\mathcal C}$ avec $\mathcal C:=\{(x,t):t>0,\ |x|<ct\}$. Montrer que $\square E=\delta_{(0,0)}$ dans $\D'(\R^2)$.
  \item Soit $\psi\in\D(\R^2)$. Expliciter $u:=\psi\star E$, montrer que $\square u=\psi$ et que $u(x,t)=0$ si $t\le\min\{s:\ (y,s)\in\supp\psi\}$.
\end{enumerate}
\end{exo}

% =====================================================================
\begin{exo}{Exercice 14 -- Vrai ou faux ?}
Pour chaque affirmation, dire si elle est vraie ou fausse, et justifier par une preuve courte ou un contre-exemple. (Viser deux à trois minutes par question.)
\begin{enumerate}[label=\arabic*.]
  \item Si $(\varphi_n)$ est une suite de $\D(\R)$ telle que $\varphi_n^{(j)}\to0$ uniformément sur $\R$ pour tout $j\in\N$, alors $\varphi_n\to0$ dans $\D(\R)$.
  \item Si $T\in\D'(\R)$ et si $\varphi\in\D(\R)$ est nulle sur $\supp T$, alors $\langle T,\varphi\rangle=0$.
  \item Si $(f_n)$ est une suite de $L^1_{\mathrm{loc}}(\R)$ qui converge presque partout vers $f\in L^1_{\mathrm{loc}}(\R)$, alors $T_{f_n}\wk T_f$.
  \item Si $T_{f_n}\wk0$ dans $\D'(\R)$, alors une sous-suite de $(f_n)$ converge vers $0$ presque partout.
  \item $\delta_0'$ est une mesure de Radon.
  \item Si $T\in\D'(\R)$ vérifie $xT=0$, alors $T=0$.
  \item Si $f\in\Cinf(\R)$ ne s'annule pas et si $fT=0$, alors $T=0$.
  \item $(x\delta_0')'=-\delta_0'$.
  \item On peut prolonger le produit $(f,T)\mapsto fT$ en donnant un sens à $h\,\delta_0$ par continuité : si $T_{g_n}\wk\delta_0$ avec $g_n\in\D(\R)$, la suite $T_{hg_n}$ a une limite indépendante du choix de $(g_n)$.
  \item Les solutions dans $\D'(\R)$ de $T'=T$ sont les $Ce^x$, $C\in\C$.
  \item Si $\psi\in\D(\R)$ est non nulle et $\psi\star T=0$, alors $T=0$.
  \item Si $f\in L^1_{\mathrm{loc}}(\R)$ et si $\langle T_f',\varphi\rangle=0$ pour toute $\varphi\in\D(\R\setminus\{0\})$, alors $f$ est constante presque partout.
  \item Si $T_n'\wk0$ dans $\D'(\R)$, alors $(T_n)$ converge dans $\D'(\R)$.
  \item La dérivée de $\ln|x|$ dans $\D'(\R)$ est la fonction $x\mapsto1/x$.
  \item Si $T\in\D'(\R)$ est d'ordre $m$, alors $T'$ est d'ordre exactement $m+1$.
  \item Si $T\in\D'(\R)$ et si $T'$ est d'ordre $0$, alors $T$ est une fonction, et même $T\in L^\infty_{\mathrm{loc}}(\R)$.
\end{enumerate}
\end{exo}

\end{document}
